\documentclass[11pt,a4paper]{article}

% -----------------------------------------------------
% PACKAGES
% -----------------------------------------------------

\usepackage[margin=1in]{geometry}
\usepackage{hyperref}
\usepackage{enumitem}
\usepackage{graphicx}
\usepackage{xcolor}
\usepackage{titling}

% -----------------------------------------------------
% TIET UCS503P PROPOSAL FORMAT
% -----------------------------------------------------

\makeatletter
\newcommand{\BvrSetLabInstructor}[1]{%
  \gdef\Bvr@LabInstructor{#1}}
\newcommand{\Bvr@LabInstructor}{Ms. Anushka Chaudhary}
\newcommand{\BvrLabInstructorValue}{\Bvr@LabInstructor}
\makeatother

% -----------------------------------------------------
% SUBTITLE
% -----------------------------------------------------

\pretitle{
    \begin{center}
    \bfseries
    \fontsize{18bp}{18bp}\selectfont
}

\makeatletter
\newcommand{\BvrSubtitle}[1]{%
    \posttitle{%
        \par
        \end{center}
        \begin{center}
        \large #1
        \end{center}
        \vskip 0.5em
    }%
}
\makeatother

% -----------------------------------------------------
% SECTION HINTS
% -----------------------------------------------------

\newcommand{\BvrHint}[1]{%
    {\normalfont\normalsize\itshape\hfill #1}
}

% -----------------------------------------------------
% PROJECT INFORMATION
% -----------------------------------------------------

\title{Sports Resource Management System}

\BvrSetLabInstructor{Ms. Anushka Chaudhary}

\BvrSubtitle{%
    Project Proposal for UCS503P  \\
    Submitted to: \BvrLabInstructorValue \\ \bfseries
    Thapar Institute of Engineering and Technology
}

\author{
    Soumya Smita Sahoo, AIML%
    \thanks{
        Roll No: \texttt{1024240041}
    }
    \and
    Himasri Chadalavada, AIML%
    \thanks{
        Roll No: \texttt{1024240038}
    }
    \and
    Gaganpreet Kaur, AIML%
    \thanks{
        Roll No: \texttt{1024240044}
    }
}

\date{September 2026}

% -----------------------------------------------------
% DOCUMENT
% -----------------------------------------------------

\begin{document}

\maketitle

% =====================================================
% 1. HIGHER-ORDER GOAL
% =====================================================

\section{Higher-order goal%
    \BvrHint{\ldots efficient campus sports management}
}

This project aims to improve the accessibility, organization, and
management of sports resources within a university campus. Sports
facilities and equipment may be available, but students, coaches, and
sports administration may face difficulty in knowing their availability,
requesting equipment, booking facilities, tracking resources, and
communicating requirements.

The immediate engineering objective is to translate these fragmented
activities into a centralized and trackable mobile application that
supports equipment management, facility booking, equipment requests,
and resource tracking.

The application is designed for different categories of users including
students or players, coaches, sports officers, administrators, and
maintenance staff.

% =====================================================
% 2. TIME-TO-VALUE
% =====================================================

\section{Time-to-value%
    \BvrHint{\ldots primary heuristic}
}

\begin{itemize}[leftmargin=0.7cm]

    \item We will deliver meaningful increments quickly by first
    implementing the core mobile workflows for login, equipment
    management, facility management, booking, and equipment requests.

    \item Development will be carried out incrementally so that the
    basic application and backend can be integrated before adding
    secondary modules.

    \item The initial prototype will focus on the most important
    equipment and facility management operations rather than attempting
    to implement every feature simultaneously.

    \item Additional functionality such as administrator management,
    maintenance tracking, feedback, and notifications will be added in
    subsequent iterations.

    \item Success in the initial stage is defined by having a working
    mobile prototype with reliable communication between the application,
    backend, and database.

\end{itemize}

% =====================================================
% 3. PROBLEM STATEMENT
% =====================================================

\section{Problem Statement}

Students and athletes often depend on sports offices or informal
communication to access campus sports equipment and facilities. This
can lead to uncertainty regarding availability, booking conflicts,
difficulty in tracking equipment, and delays in handling resource
requests.

Common pain points include:

\begin{itemize}[leftmargin=0.7cm]

    \item \textbf{Fragmentation:}
    Equipment requests, facility bookings, and resource-related
    communication may be handled through different channels such as
    physical registers, messages, calls, or informal communication.

    \item \textbf{Low availability transparency:}
    Students and coaches may not know whether a particular equipment
    item or facility is currently available.

    \item \textbf{Booking conflicts:}
    Multiple users may attempt to book the same facility for overlapping
    time periods, requiring manual coordination.

    \item \textbf{Poor equipment tracking:}
    Sports administration may face difficulty tracking equipment
    quantities, availability, condition, requests, and usage.

    \item \textbf{Equipment identification:}
    When equipment requires repair, maintenance staff need to identify
    the exact item. A unique Equipment ID can be used for this purpose.

    \item \textbf{Request management:}
    Equipment requests need to be stored and managed systematically so
    that sports staff can review and process them.

    \item \textbf{Limited centralized information:}
    Users, equipment, facilities, bookings, and requests should be
    connected through a centralized database.

\end{itemize}

% =====================================================
% 4. PROPOSED SOLUTION
% =====================================================

\section{Proposed Solution%
    \BvrHint{\ldots Sports Resource Management Mobile Application}
}

\subsection{Overview}

Build a cross-platform mobile application that provides role-based
access to sports resource management functionality.

The proposed system contains the following components:

\begin{itemize}[leftmargin=0.7cm]

    \item \textbf{Student/Player App:}
    Students can log in, view sports equipment, check availability,
    request equipment, view facilities, and book available facilities.

    \item \textbf{Coach Module:}
    Coaches can view facilities, check availability, book facilities
    for sports activities, and request required equipment.

    \item \textbf{Sports Officer Module:}
    Sports officers can manage equipment and facilities, review
    equipment requests, and manage booking-related information.

    \item \textbf{Maintenance Module:}
    Maintenance staff will be able to identify equipment using its
    Equipment ID and manage repair-related information. This module is
    currently under development.

    \item \textbf{Administrator Module:}
    The administrator role is included for system-level management such
    as user and role management. This functionality is currently under
    development.

    \item \textbf{Feedback and Notification Modules:}
    Feedback and notifications are part of the planned system design.
    These modules are currently under development.

\end{itemize}

The system is being developed as a mobile application only. There is no
separate web-based administrator dashboard in the current project scope.

\subsection{Core workflow}

\subsubsection*{Equipment workflow}

\begin{enumerate}[leftmargin=0.7cm]

    \item The student or coach logs into the mobile application.

    \item The user views the available sports equipment.

    \item The application displays relevant equipment information,
    including availability and status.

    \item The user selects the required equipment and submits a request.

    \item The Sports Officer reviews the submitted request.

    \item If the request is approved, the equipment can be issued to the
    user.

    \item The equipment record can be updated with relevant issue and
    return information.

    \item If the equipment requires repair, its unique Equipment ID can
    be used to associate it with a maintenance record.

\end{enumerate}

\subsubsection*{Facility booking workflow}

\begin{enumerate}[leftmargin=0.7cm]

    \item The student or coach selects a sports facility.

    \item The application displays information about the selected
    facility and its availability.

    \item The user selects the required date and time slot.

    \item The system checks the existing booking information to
    determine whether the selected slot is available.

    \item If the slot is available, the booking information is stored
    in the database.

    \item The user can subsequently access the booking information
    through the application.

\end{enumerate}

\subsubsection*{Maintenance workflow}

The maintenance process begins when an equipment item is reported as
requiring attention. Maintenance staff first inspect the equipment to
understand the problem. If repair is required, the equipment is marked
as being under repair while the maintenance work is carried out. After
the repair is completed, the equipment is inspected again and its status
is updated. If the equipment is found to be usable, it is marked as
available for future use.

Each maintenance record is associated with the corresponding Equipment
ID so that maintenance staff can identify the exact equipment under
repair.

This module is part of the planned system and is not fully functional
in the current prototype.

\subsection{Operational constraints%
    \BvrHint{\ldots campus mobile setting}
}

\begin{itemize}[leftmargin=0.7cm]

    \item The application should use a lightweight and responsive
    mobile interface suitable for commonly available student
    smartphones.

    \item Core functions such as equipment viewing, booking, and
    equipment requests should require minimal steps.

    \item The application should provide clear loading and error states
    when communication with the backend is delayed or unavailable.

    \item The system should prevent conflicting facility bookings.

    \item Equipment marked as unavailable or under maintenance should not
    be treated as normally available.

    \item The system will prioritize workflow correctness and reliable
    resource management rather than advanced prediction algorithms.

    \item The backend and database will remain separate from the mobile
    interface so that additional modules can be integrated later.

\end{itemize}

% =====================================================
% 5. SOLUTION APPROACH
% =====================================================

\section{Solution Approach%
    \BvrHint{\ldots engineering focus}
}

This is an engineering course project, so the primary focus is on
mobile workflow design, role-based access, resource management, database
integration, system correctness, and iterative development rather than
novel algorithmic research.

\subsection{Resource availability and conflict checking}

Facility and equipment availability will be determined using stored
resource information, booking status, equipment quantity, condition,
and maintenance status.

The backend will validate relevant availability information before
storing booking-related records. This is intended to reduce conflicting
facility reservations.

\subsection{Equipment lifecycle management}

Equipment records will contain a unique Equipment ID along with
information such as name, category, quantity, condition, and status.

The equipment lifecycle begins when an item is available for use. When
a user receives the equipment, its status can be recorded as issued.
After the equipment is returned, its condition can be checked. The
equipment can then either be made available again or associated with a
maintenance process if repair is required.

The use of Equipment IDs also allows individual pieces of equipment to
be distinguished when multiple similar resources are present.

\subsection{Role-based access}

The application begins with role selection followed by the login
workflow.

The selected role determines the functionality available to the user.

The main roles are:

\begin{itemize}[leftmargin=0.7cm]

    \item Student / Player
    \item Coach
    \item Sports Officer
    \item Administrator
    \item Maintenance Staff

\end{itemize}

\subsection{Mobile application architecture}

The system uses a mobile client-server architecture.

The Flutter and Dart mobile application provides the interface through
which users interact with the system. Requests from the application are
handled by the Node.js backend, which contains the server-side
application logic. The backend communicates with PostgreSQL whenever
data needs to be stored or retrieved.

This separation allows the mobile application, backend, and database to
be developed and maintained as separate components.

\subsection{Technology stack}

The project uses the following technologies:

\begin{itemize}[leftmargin=0.7cm]

    \item \textbf{Flutter:}
    Cross-platform mobile application framework.

    \item \textbf{Dart:}
    Programming language used for the mobile application.

    \item \textbf{Node.js:}
    Backend runtime and server-side implementation.

    \item \textbf{PostgreSQL:}
    Relational database for persistent application data.

    \item \textbf{Visual Studio Code:}
    Development environment.

    \item \textbf{Git and GitHub:}
    Source-code version control and collaboration.

\end{itemize}

% =====================================================
% 6. EVALUATION CRITERION
% =====================================================

\section{Evaluation Criterion%
    \BvrHint{\ldots measurable after implementation}
}

The current submission is a prototype-stage project. Therefore, formal
evaluation results are not claimed at this stage.

After the remaining modules have been implemented, the system can be
evaluated using functional and usability-oriented metrics.

\subsection{Primary evaluation criteria}

\textbf{Core workflow correctness:}

The main workflows should operate correctly from the mobile interface
through the backend and database.

The following should be verified:

\begin{itemize}[leftmargin=0.7cm]

    \item Successful role selection and login.

    \item Correct display of equipment information.

    \item Correct submission and storage of equipment requests.

    \item Correct display of facilities.

    \item Correct facility booking.

    \item Correct handling of booking conflicts.

    \item Correct storage and retrieval of booking information.

\end{itemize}

\subsection{Secondary evaluation criteria}

\begin{itemize}[leftmargin=0.7cm]

    \item \textbf{Availability accuracy:}
    Equipment and facility availability shown by the application should
    correspond to the database records.

    \item \textbf{Booking correctness:}
    The system should not confirm conflicting bookings.

    \item \textbf{Request correctness:}
    Equipment requests should be stored with the correct user and
    request information.

    \item \textbf{Equipment identification:}
    Equipment IDs should uniquely identify equipment records.

    \item \textbf{Usability:}
    Common workflows such as viewing resources, making a booking, and
    submitting a request should be understandable to users.

    \item \textbf{Reliability:}
    Core application functionality should remain usable during normal
    prototype testing.

\end{itemize}

\subsection{Future validation plan}

Once the complete prototype is available, representative users can be
asked to perform common tasks such as:

\begin{itemize}[leftmargin=0.7cm]

    \item Log in using their assigned role.

    \item Search or view equipment.

    \item Submit an equipment request.

    \item View facility availability.

    \item Book a facility.

    \item Attempt a conflicting booking.

    \item Review booking information.

    \item Test equipment maintenance identification using Equipment ID.

\end{itemize}

The results of such testing will be documented in a later project
iteration.

% =====================================================
% 7. SCALABILITY
% =====================================================

\section{Scalability%
    \BvrHint{\ldots with foresight}
}

The system should support growth in the number of students, coaches,
equipment items, facilities, bookings, and requests.

\begin{enumerate}[leftmargin=0.7cm]

    \item \textbf{Scalable architecture:}
    The mobile client, backend, and database are separated so that each
    layer can be developed and optimized independently.

    \item \textbf{Database scalability:}
    Database indexing can be introduced for frequently searched
    equipment, facility, and booking information.

    \item \textbf{Large resource lists:}
    Pagination can be introduced when the equipment and booking
    catalogue becomes large.

    \item \textbf{Efficient search:}
    Search and filtering can be improved for large campus inventories.

    \item \textbf{Backend scalability:}
    The Node.js backend can be extended to support a larger number of
    concurrent mobile users.

    \item \textbf{Future deployment:}
    The architecture can later be deployed using cloud-hosted backend
    and managed database services.

\end{enumerate}

% =====================================================
% 8. ENGINE AVAILABILITY HEURISTIC
% =====================================================

\section{Engine availability heuristic}

A mature set of development tools and technologies is available to
implement the proposed mobile application:

\begin{itemize}[leftmargin=0.7cm]

    \item \textbf{Flutter} for cross-platform mobile development.

    \item \textbf{Dart} for application programming.

    \item \textbf{Node.js} for backend API development.

    \item \textbf{PostgreSQL} for structured relational data.

    \item \textbf{Git/GitHub} for version control and collaboration.

    \item Authentication mechanisms for role-based access.

    \item Testing tools for validating application and backend
    functionality.

\end{itemize}

These mature technologies allow the project to focus primarily on
workflow design, resource management, database integration, and system
correctness rather than developing new infrastructure.

% =====================================================
% 9. PROJECT SCOPE AND DELIVERABLES
% =====================================================

\section{Project scope and deliverables%
    \BvrHint{\ldots iteration-friendly}
}

\subsection{Initial deliverable%
    \BvrHint{\ldots working prototype}
}

The initial prototype focuses on the following functionality:

\begin{itemize}[leftmargin=0.7cm]

    \item Flutter mobile application setup.

    \item Role selection.

    \item User login.

    \item Student / Player functionality.

    \item Coach functionality.

    \item Sports Officer functionality.

    \item Equipment catalogue and equipment information.

    \item Equipment availability.

    \item Equipment request submission.

    \item Facility listing.

    \item Facility availability.

    \item Basic facility booking.

    \item Node.js backend.

    \item PostgreSQL database integration.

    \item Backend and mobile application communication.

\end{itemize}

\subsection{Subsequent deliverables}

The following modules will be developed incrementally:

\begin{itemize}[leftmargin=0.7cm]

    \item Administrator functionality.

    \item Maintenance Staff interface.

    \item Equipment-specific maintenance tracking.

    \item Equipment damage and repair status tracking.

    \item Improved equipment issue and return tracking.

    \item Feedback and suggestion functionality.

    \item Notification functionality.

    \item Booking and equipment history.

    \item Improved search and filtering.

    \item Additional resource management features.

    \item Comprehensive testing and evaluation.

\end{itemize}

% =====================================================
% 10. RISKS AND MITIGATIONS
% =====================================================

\section{Risks and mitigations}

\begin{itemize}[leftmargin=0.7cm]

    \item \textbf{Incorrect inventory information:}
    Maintain centralized equipment records and ensure that equipment
    status is updated through the application.

    \item \textbf{Double booking:}
    Perform availability validation before confirming a facility
    booking.

    \item \textbf{Equipment identification problems:}
    Assign a unique Equipment ID to each trackable equipment record.

    \item \textbf{Equipment condition disputes:}
    Maintain condition information when equipment is issued and
    returned. A more detailed condition workflow can be introduced
    later.

    \item \textbf{Incomplete prototype modules:}
    Prioritize core functionality first and implement administrator,
    maintenance, feedback, and notification modules incrementally.

    \item \textbf{Network/connectivity issues:}
    Provide appropriate loading, retry, timeout, and error states in
    the mobile application.

    \item \textbf{Device compatibility:}
    Test the mobile application on representative devices and screen
    sizes.

    \item \textbf{User adoption:}
    Keep common workflows such as login, resource viewing, booking, and
    equipment requests simple.

\end{itemize}

% =====================================================
% 11. SUMMARY
% =====================================================

\section{Summary}

This project proposes a centralized mobile Sports Resource Management
System for managing university sports equipment and facilities.

The application connects students, coaches, sports officers,
administrators, and maintenance staff through role-based mobile
workflows. Core functionality includes equipment management, facility
availability, facility booking, and equipment requests.

The system is being developed using Flutter and Dart for the mobile
application, Node.js for the backend, and PostgreSQL for the database.
The separation between the mobile client, backend, and database provides
a foundation for future expansion.

The current prototype focuses on the core working functionality.
Administrator functionality, maintenance tracking, feedback, and
notifications are planned as subsequent modules rather than being
claimed as completed features.

The use of unique Equipment IDs will allow individual equipment items to
be identified when they are issued, returned, damaged, or eventually
sent for maintenance.

The project therefore focuses on practical software engineering
requirements including multi-user workflows, role-based access,
database design, mobile application development, backend integration,
system correctness, iterative development, and future scalability.

\end{document}
